\documentclass[11pt, a4paper]{article}

\usepackage[T1]{fontenc}
\usepackage[utf8]{inputenc}
\usepackage{tgpagella}
\usepackage[spanish, es-noshorthands]{babel}
\usepackage{amsmath, amssymb}
\usepackage{booktabs}
\usepackage{tabularx}
\usepackage[most]{tcolorbox}
\usepackage[margin=2.5cm]{geometry}
\usepackage{fancyhdr}

\pagestyle{fancy}
\fancyhf{}
\setlength{\headheight}{16pt}
\lhead{\textbf{Ing. Mecatrónica}}
\chead{Circuitos Electrónicos II (IMT-221)}
\rhead{Semana 7 - Sesión 2}
\lfoot{Pre-Lab 3: Errores DC}
\rfoot{Página \thepage}
\renewcommand{\headrulewidth}{0.4pt}
\definecolor{ucbblue}{RGB}{0, 51, 102}

\begin{document}

\begin{center}
    \Large\textbf{\textcolor{ucbblue}{Guía de Pre-Lab 3: Diagnóstico y Compensación DC}}[cite: 11]
\end{center}

\begin{tcolorbox}[colback=red!5, colframe=red!70!black, title=\textbf{ADVERTENCIA IMPORTANTE}]
    \textbf{PRE-LAB 3 — NO REEMPLAZA LA EJECUCIÓN FÍSICA DEL LABORATORIO.[cite: 11]} La actividad queda incompleta hasta implementar, medir, diagnosticar, compensar, volver a medir e interpretar el circuito físico en la Semana 8[cite: 11].
\end{tcolorbox}

\section*{A. Definición del Circuito y Datasheet[cite: 11]}
\textbf{Op-Amp Oficial:} LM741, encapsulado DIP-8[cite: 11]. \\
\textbf{Fabricante físico:} \rule{4cm}{0.4pt} \textbf{Documento / revisión:} \rule{4cm}{0.4pt}[cite: 11] \\
\textbf{Alimentación de laboratorio:} $V^+ = +15\,\text{V}$, $V^- = -15\,\text{V}$[cite: 11]. \\
\textbf{Red diferencial ($1\%$):} $R_1=R_3=10\,\text{k}\Omega$, $R_2=R_4=100\,\text{k}\Omega$[cite: 11]. \\
\textbf{Elemento de compensación:} Trimmer/potenciómetro lineal de $10\,\text{k}\Omega$[cite: 11]. \\
\textbf{Desacoplo local mínimo:} $2 \times 100\,\text{nF}$ cerámicos[cite: 11]. \\
\textbf{Pinout verificado:} Sí / No \hfill \textbf{Pines OFFSET NULL:} \rule{4cm}{0.4pt}[cite: 11] \\
\textbf{Input offset voltage ($V_{OS}$):} \rule{3cm}{0.4pt} \textbf{Condición de prueba:} \rule{3cm}{0.4pt}[cite: 11] \\
\textbf{Input bias current ($I_B$):} \rule{3cm}{0.4pt} \textbf{Condición de prueba:} \rule{3cm}{0.4pt}[cite: 11] \\
\textbf{Input offset current ($I_{OS}$):} \rule{3cm}{0.4pt} \textbf{Procedimiento OFFSET NULL:} \rule{3cm}{0.4pt}[cite: 11]

\section*{B. Análisis Ideal[cite: 11]}
Derive simbólicamente, para el circuito confirmado: $V_{o,\mathrm{ideal}} = f(V_1, V_2, R_1, R_2, R_3, R_4)$[cite: 11].
Verifique matemáticamente si se cumple $\frac{R_2}{R_1} = \frac{R_4}{R_3}$[cite: 11]. Dado que la razón es 10, confirme la ecuación $V_{o,\mathrm{ideal}} = 10(V_2 - V_1)$[cite: 11].

\section*{C. Condiciones de Prueba Planificadas (DC)[cite: 11]}
\begin{table}[h]
    \centering
    \renewcommand{\arraystretch}{1.3}
    \begin{tabularx}{\textwidth}{|c|X|X|X|}
    \hline
    \textbf{Condición} & \textbf{$V_1$ (V)} & \textbf{$V_2$ (V)} & \textbf{Propósito} \\ \hline
    T1 & \rule{3cm}{0.4pt} & \rule{3cm}{0.4pt} & Baseline diferencial[cite: 11] \\ \hline
    T2 & \rule{3cm}{0.4pt} & \rule{3cm}{0.4pt} & Condición de salida ideal conocida[cite: 11] \\ \hline
    T3 & \rule{3cm}{0.4pt} & \rule{3cm}{0.4pt} & Diagnóstico / Compensación[cite: 11] \\ \hline
    \end{tabularx}
\end{table}

\section*{D. Prediction Table (Columnas físicas vacías hasta Semana 8)[cite: 11]}
\begin{table}[h]
    \centering
    \small
    \renewcommand{\arraystretch}{1.5}
    \begin{tabularx}{\textwidth}{|c|c|c|c|c|X|}
    \hline
    \textbf{Test} & \textbf{Ideal $V_o$} & \textbf{LTSpice $V_o$} & \textbf{Physical $V_o$ (Before)} & \textbf{Physical $V_o$ (After)} & \textbf{Interpretation} \\ \hline
    T1 & & & & & \\ \hline
    T2 & & & & & \\ \hline
    T3 & & & & & \\ \hline
    \end{tabularx}
\end{table}

\section*{E. Nodos Requeridos de Medición[cite: 11]}
Identificar: $V_1, V_2, V_+, V_-, V_o$, Alimentación $+$, Alimentación $-$, Nodos de compensación[cite: 11].

\section*{F. Procedimiento Baseline (Para Semana 8)[cite: 11]}
1. Verificar circuito sin alimentación. 2. Medir valores reales resistivos. 3. Verificar orientación y pinout Op-Amp. 4. Establecer alimentación. 5. Aplicar $V_1, V_2$. 6. Medir $V_1, V_2$. 7. Medir $V_o$. 8. Registrar error ($e = V_{o,\mathrm{fis}} - V_{o,\mathrm{ideal}}$). 9. \textbf{NO compensar todavía.} 10. Diagnosticar[cite: 11].

\section*{G. Preguntas de Diagnóstico y Plan de Compensación[cite: 11]}
¿Las fuentes, tierra y pinout son correctos? ¿El error persiste con entradas balanceadas? ¿El instrumento mide respecto al nodo correcto?[cite: 11] \\
\textbf{Método previsto:} \rule{6cm}{0.4pt} \textbf{Variable a ajustar:} \rule{4cm}{0.4pt}[cite: 11] \\
\textbf{Condición durante compensación:} \rule{4cm}{0.4pt} \textbf{Criterio de mejora:} \rule{3cm}{0.4pt}[cite: 11] \\

\section*{H. Hipótesis y Evidencia[cite: 11]}
Antes del laboratorio: ¿Qué desviación respecto al modelo ideal espera y qué causa considera plausible?[cite: 11] \\
\rule{\textwidth}{0.4pt} \\
Requisito Sem 8: Fotos del montaje, medición baseline, error, diagnóstico y comparación before/after[cite: 11].

\end{document}
